\documentclass[11pt]{article}
\usepackage[margin=1in]{geometry}
\usepackage{amsmath,amssymb}
\usepackage{microtype}
\usepackage{xcolor}
\usepackage{hyperref}
\hypersetup{colorlinks=true,urlcolor=blue!55!black,citecolor=blue!55!black}
\title{Literature status of the infinite-dimensional direction\\from arXiv:2311.09207}
\author{Scoped literature-identification packet}
\date{11 August 2026}

\begin{document}
\maketitle

\begin{center}
\fbox{\parbox{0.91\linewidth}{\textbf{Status:} literature-implied answer to the broad infinite-dimensional direction.  The later paper does not separately restate the source's exact quasi-Metropolis time-domain formula, so that formula-specific scope is not claimed here.}}
\end{center}

\section{Source direction}

Chen, Kastoryano, and Gily\'en construct an efficient exactly
detailed-balanced noncommutative Gibbs sampler in finite dimension~\cite{CKG}.
After Proposition B.2, the source states on PDF page 32 that its exact coherent
term appears to persist in infinite dimension and leaves verification of the
construction and analysis for future work.

The mathematically meaningful extension must handle an unbounded Hamiltonian:
on an infinite-dimensional Hilbert space a bounded Hamiltonian cannot have a
normal Gibbs density, because $e^{-\beta H}$ is bounded below by a positive
multiple of the identity and therefore is not trace class.

\section{Later result and identification}

Becker, Rouz\'e, and Salzmann explicitly cite the source as the first exact
Lindbladian Gibbs sampler and then develop a rigorous infinite-dimensional
theory for unbounded Hamiltonians on separable Hilbert spaces~\cite{BRS}.  Their
abstract and introduction state that they construct KMS-symmetric quantum
Markov semigroups that are well posed and implementable on qubit hardware.

More precisely, Proposition 2.8 identifies a trace-class generator under
explicit compatibility and boundedness hypotheses.  Section 4.1 introduces
the Gaussian-convoluted generator
\[
 \mathcal L_{\sigma_E,\widehat f,H}
\]
with the same characteristic coherent coefficient
\[
 \frac{i}{2}\tanh\!\left(\frac{\beta(\nu_1-\nu_2)}4\right)
 (A_{\nu_1})^*A_{\nu_2},
\]
and Proposition 4.3 proves that its closed form generates a completely Markov
KMS-symmetric semigroup.  The paper further proves controlled finite-rank
approximations and circuit implementations, including singular
Metropolis-type filters.

This supplies the rigorous generation, exact KMS symmetry, and truncation
analysis requested by the broad future-work direction.  The identification is
classified as \emph{literature implied}, rather than
\emph{literature already answered}, because the later authors do not advertise
their theorem as a line-by-line verification of Proposition B.2 in
arXiv:2311.09207.

\section{Scope limitation}

The later construction uses an equivalent subsequent sampler framework and
Dirichlet-form generation theory.  This packet therefore does not assert that
arXiv:2604.01192 proves the validity of the exact quasi-Metropolis time-domain
formula in Proposition B.2 under every conceivable domain hypothesis.  It
does establish that the broad infinite-dimensional exact-detailed-balance
direction is no longer an open new-result target.

\section{Search evidence}

The bounded search through 11 August 2026 covered the run indexes; exact
phrases from the source's future-work sentence; citations to arXiv:2311.09207;
the finite-dimensional follow-up arXiv:2404.05998; and the 2026
infinite-dimensional papers arXiv:2604.01192, 2604.06077, and 2604.15263.  The
decisive source is arXiv:2604.01192; arXiv:2604.06077 explicitly invokes it as
the companion infinite-dimensional framework.

\begin{thebibliography}{9}
\bibitem{CKG}
C.-F. Chen, M.~J. Kastoryano, and A. Gily\'en,
\emph{An efficient and exact noncommutative quantum Gibbs sampler},
arXiv:2311.09207.
\url{https://arxiv.org/abs/2311.09207}.

\bibitem{BRS}
S. Becker, C. Rouz\'e, and R. Salzmann,
\emph{Quantum Gibbs Sampling in Infinite Dimensions: Generation, Mixing Times
and Circuit Implementation}, arXiv:2604.01192 (2026).
\url{https://arxiv.org/abs/2604.01192}.
\end{thebibliography}

\end{document}
